% 02_body.tex — 산출물 2/5 프로젝트 헌장 (빈 템플릿 / 완성 예시 공용 본문)
% 드라이버: 02_프로젝트헌장_빈템플릿.tex (\wbblanktrue) / 02_프로젝트헌장_완성예시.tex (\wbblankfalse)
% 내용 출처: PM트랙/Day1_워크북.md §2.3(항목 가이드) §2.4(빈 템플릿) §2.5(온담 P1 완성 예시본)
\documentclass[10pt,a4paper]{article}
\usepackage[a4paper,margin=16mm,top=14mm,bottom=20mm]{geometry}
\def\DocNum{2}
\def\DocTitle{프로젝트 헌장}
\def\DocSub{Project Charter}
\def\DocVariant{\B{빈 템플릿}{온담 P1 완성 예시}}
\usepackage{wbstyle}
\begin{document}

\docheader
 {\Bg{[정식 명칭 / 코드]}{차세대 주문·재고 통합 플랫폼 구축 / ONE ONDAM P1}}
 {\Bg{[v1.x / 작성일]}{v1.0 / 2026-01-05 (G0)}}
 {\Bg{[PM 직책·실명]}{P1 프로젝트 매니저 (발주사 IT본부)\rd{7mm}}}
 {\Bg{[스폰서 직책·실명]}{정미란 재무본부장 CFO (프로그램 스폰서·Executive)\rd{7mm}}}

% ------------------------------------------------------------------ 1
\secbar{1. 문서 정보와 승인}
\guide{정식 명칭·코드, 상위 프로그램, 헌장 버전, 작성자(PM), 승인자(스폰서 — \textbf{개인의 직책과 실명}, 위원회는 서명하지 않는다)와 승인일. 프로그램 하위 프로젝트면 프로그램 헌장과의 관계를 명시.}
{\footnotesize
\begin{tabularx}{\linewidth}{|L{28mm}|Y|L{22mm}|Y|}\hline
\hd{항목} & \hd{내용} & \hd{항목} & \hd{내용} \\ \hline
\B{\rd{9mm}}{}프로젝트명 / 코드 & \Bg{[정식 명칭 / 코드]}{차세대 주문·재고 통합 플랫폼 구축 / ONE ONDAM P1} & 상위 프로그램 & \Bg{[프로그램명, 프로그램 헌장과의 관계]}{ONE ONDAM (프로그램 헌장 2026-01-05 G0 서명본의 하위 문서)} \\ \hline
\B{\rd{9mm}}{}버전 / 작성일 & \Bg{[v1.x / 날짜]}{v1.0 / 2026-01-05} & 작성자(PM) & \Bg{[직책·실명]}{P1 프로젝트 매니저 (발주사 IT본부 소속)} \\ \hline
\B{\rd{9mm}}{}승인자(스폰서) & \Bg{[직책·실명]}{정미란 재무본부장 CFO (프로그램 스폰서·Executive)} & 승인일 & \Bg{[날짜 / 게이트]}{2026-01-05 (G0)} \\ \hline
\end{tabularx}}

% ------------------------------------------------------------------ 2
\secbar{2. 프로젝트 목적과 사업적 정당성 (Purpose \& Justification)}
\guide{3~5문장. 문제의 현재 측정값 → 기대 효과의 목표값과 금액 → 인용 문서(비즈니스 케이스 요약). 숫자를 반드시 포함. 솔루션을 목적으로 쓰지 말 것(“플랫폼을 구축한다”는 수단이다). 프로그램 전체의 목적을 그대로 옮기지 말 것.}
\begin{fieldbox}
\ifwbblank
\lines{7}{9.5mm}
\else
온담F\&B홀딩스는 615개 매장(직영 227·가맹 388), 이천 물류센터, 안성·이천 공장의 주문과 재고를 여섯 개의 분리된 시스템으로 관리하고 있으며, 그 결과 폐기율 3.9\%(식자재 취급원가 2,530억 기준 연 98.7억 손실), 결품률 4.7\%, 재고 정확도 91.3\%, 가맹 발주 리드타임 38시간의 운영 손실을 안고 있다. 매장 시스템 OD-POS 2.1은 서버 OS 지원이 2024-07에 종료되어 보안 패치 없이 결제를 처리하고 있고, 개인정보 보호법 개정 대응(2026-09-11)과 전자금융업 등록 심사(2027-03-31)를 현행 아키텍처로는 충족할 수 없다. 본 프로젝트는 매장·물류·공장의 주문과 재고를 하나의 원장으로 실시간 관리하는 플랫폼을 구축해 M+24(2028-01) 기준 연간 직접 편익 118억(폐기율 43·배달수수료 회피 37·매장 인건비율 38)을 실현하고, 규제 기한 안에 OD-POS를 대체하며, 2027년 3분기 IPO 예비심사 전에 플랫폼을 안정 가동하는 것을 목적으로 한다. 사업적 정당성의 상세는 「비즈니스 케이스 요약 v0.9」(산출물 1)에 따르며, 편익 실현률 손익분기는 66\%다.
\fi
\end{fieldbox}

% ------------------------------------------------------------------ 3
\secbar[80mm]{3. 측정 가능한 목표와 성공 기준 (Objectives \& Success Criteria)}
\guide{두 층으로 나눈다 — \textbf{프로젝트 관리 성공}(기한·예산·범위·품질 안에 끝났는가, 종료 게이트 판정)과 \textbf{프로젝트 성공}(편익이 실현되었는가, 종료 후 M+n 판정). 각 행이 참/거짓으로 판정 가능해야 하며 허용범위를 목표값에 포함한다. 판정자와 판정 시점 없는 행은 성공 기준이 아니다.}
{\footnotesize
\begin{tabularx}{\linewidth}{|L{17mm}|L{24mm}|Y|Y|L{18mm}|L{26mm}|}\hline
\hd{구분} & \hd{목표} & \hd{측정 지표} & \hd{목표값(허용범위)} & \hd{판정 시점} & \hd{판정자} \\ \hline
\ifwbblank
\rd{9mm}관리 성공 & \g{[기한 내 종료 등]} & \g{[무엇을 재는가]} & \g{[값 (허용범위)]} & \g{[종료 게이트]} & \g{[직책·실명]} \\ \hline
\rd{9mm}관리 성공 & & & & & \\ \hline
\rd{9mm}관리 성공 & & & & & \\ \hline
\rd{9mm}관리 성공 & & & & & \\ \hline
\rd{9mm}관리 성공 & & & & & \\ \hline
\rd{9mm}프로젝트 성공 (편익) & & & & \g{[종료 후 M+n]} & \\ \hline
\rd{9mm}프로젝트 성공 (편익) & & & & & \\ \hline
\rd{9mm}프로젝트 성공 (편익) & & & & & \\ \hline
\rd{9mm}프로젝트 성공 (편익) & & & & & \\ \hline
\else
관리 성공 & 기한 내 운영 이관 & G4 완료일 & 2027-05-29 (+15영업일 이내) & G4 & 정미란 스폰서 \\ \hline
관리 성공 & 집행 한도 내 완료 & 실제 원가 & 156.0억 이내 (+6\% = 10.08억 이내 허용) & G4 & 재경팀장 검증, 스폰서 판정 \\ \hline
관리 성공 & 범위 완전성 & RTM 커버리지 (요구사항 1,180건 → 시험케이스 → 검수확인) & 100\% (동결 범위 기준, ±3\% 변동 허용) & G3 & 정하늘 수석감리원 확인, 박세연 승인 \\ \hline
관리 성공 & 품질 & 결함밀도 & 0.4건/FP 이하 (UAT 종료 기준) & G3 & 수행사 PM 보고, 발주사 PMO 검증 \\ \hline
관리 성공 & 운영 인수 & SAC 미해소 항목 & 0건 또는 해소 기한·책임자 서명 & G4 & 박준영 서비스운영팀장 \\ \hline
관리 성공 & 형상 인수 & 소스·형상 발주사 저장소 인수율 & 100\% (2023 감사 지적 대응) & G4 & 박세연 IT본부장 \\ \hline
프로젝트 성공 & 가맹 발주 리드타임 & 발주 접수→출고 확정 & 38h → 12h & M+12 (2027-01) & 오정근 Senior User \\ \hline
프로젝트 성공 & 폐기율 & 폐기액 ÷ 식자재 취급원가 & 3.9\% → 2.2\% (−5\%p 편익 허용범위) & G5 2028-03-31 & 오정근 Senior User, 재무본부 판정 \\ \hline
프로젝트 성공 & 재고 정확도 & 실사 일치율 & 91.3\% → 98.5\% (분모 확장 방안 M+3 확정) & G5 & 한동석 물류본부장 \\ \hline
프로젝트 성공 & 결품률 & 매장 결품률 & 4.7\% → 1.5\% (산정식 M+6 확정) & G5 & 오정근 Senior User \\ \hline
\fi
\end{tabularx}}

% ------------------------------------------------------------------ 4
\secbar[118mm]{4. 상위 요구사항 (High-level Requirements)}
\guide{이해관계자 수준(StRS)의 요구 5~10줄 — 상세 명세(SRS)가 아니다. 줄마다 (기능 / 비기능 구분, 요구 제공자 실명)을 괄호로 적는다. 가용성·성능·보안·규제 같은 비기능 요구를 빠뜨리지 말 것.}
\begin{fieldbox}
\ifwbblank
\foreach \i in {1,...,8}{\g{\i.} \rule{0pt}{9mm}\hrulefill\ \g{(기능/비기능 · 요구 제공자: \hspace{22mm})}\par}
\else
\begin{enumerate}
\item 매장·물류센터·공장 출하 재고를 단일 원장에서 준실시간(15분 이내)으로 조회할 수 있다. (기능 / 박세연 IT본부장, 한동석 물류본부장)
\item 자사앱·배달3사·POS·가맹 발주·B2B EDI의 주문이 하나의 주문 원장에 기록되고 매장 간 재고 이동·픽업이 가능하다. (기능 / 오정근 영업본부장)
\item 가맹점주가 포털에서 발주·마감·출고 확인·배송 추적을 한 화면에서 처리할 수 있다. (기능 / 서정필 가맹점주협의회장, 오정근)
\item 성수기(7~8월·12월·명절 전 2주)에 매장 시스템 변경이 발생하지 않도록 배포 창이 설계되어 있다. (비기능·운영 / 오정근)
\item 운영 가용성 99.5\% 이상, 월 허용 다운타임 3시간 39분 이내, 컷오버 시 롤백 절차 보유. (비기능 / 박준영 서비스운영팀장)
\item 개인정보 처리항목 142개·위탁사 19개·재위탁 4건이 설계에 반영되고 망분리·전자금융 요건을 충족한다. (비기능·규제 / 최영주 CPO)
\item IT본부 47명이 벤더 없이 운영·변경할 수 있는 아키텍처와 문서(인터페이스 47본 설계문서 포함)를 갖춘다. (비기능 / 박세연)
\item SAP ECC 6.0과의 연동은 2027-12 이후 마이그레이션을 고려한 구조(애드온·API 분리)로 설계한다. (비기능 / 박세연)
\end{enumerate}
\fi
\end{fieldbox}

% ------------------------------------------------------------------ 5
\secbar[\B{195mm}{70mm}]{5. 범위 개요 — 포함과 제외 (Scope: In / Out)}
\guide{포함 5~8줄, 제외 5~8줄(\textbf{같은 무게로}), 경계 조건 3~5줄(“~의 기술 요건을 반영하되 ~자체는 하지 않는다”). 제외 항목마다 그것을 누가 하는지(프로젝트 코드/부서)를 적는다. “기타 부대 업무 일체”처럼 범위를 열어 두지 말 것.}
\begin{fieldbox}
\textbf{포함}\par
\ifwbblank
\foreach \i in {1,...,5}{\g{–} \rule{0pt}{8.5mm}\hrulefill\par}
\else
\begin{itemize}
\item 신규 클라우드 POS·매장 운영 시스템 — 직영 227점 단말 477대 신규 교체, 가맹 388점 기존 단말에 SW 배포
\item 주문 허브(Order Hub) — 전 채널 주문 원장, 옴니채널(픽업·매장 간 이동) 시나리오
\item 통합 재고 관리 — 매장·이천센터·공장 출하 재고 단일 원장, 3PL WMS 연동을 야간 배치에서 준실시간 API로 전환(연동 명세 작성 및 P1측 구현)
\item 가맹 발주 포털 — 전화·팩스·엑셀 39\% 흡수, 발주 리드타임 12시간
\item 통합 계층과 데이터 전환 — API 게이트웨이, 상품·고객 MDM, SAP 애드온, 인터페이스 47본 재구축·재연결 및 설계문서화, 1,842 테이블·7.4TB 이관
\end{itemize}
\fi
\end{fieldbox}
\begin{fieldbox}
\textbf{제외} \guide{(담당: 프로젝트 코드 / 부서)}\par
\ifwbblank
\foreach \i in {1,...,6}{\g{–} \rule{0pt}{8.5mm}\hrulefill\ \g{(담당: \hspace{25mm})}\par}
\else
\begin{itemize}
\item SAP ECC 6.0 자체의 교체 또는 업그레이드 (담당: 2027-12 유지보수 종료 대응 별도 과제, IT본부)
\item 이천 물류센터 설비·자동화 및 WMS 자체 (담당: P2)
\item 매장 직원·가맹점주 교육, 슈퍼유저 양성, 가맹 계약 부속서 개정 협상 (담당: P3, 협상 책임 오정근)
\item 개인정보 보호법 대응 완료 판정, 전자금융업 등록 신청, ISMS-P 인증 취득 자체 (담당: P4)
\item 운영 조직 편성, SLA 체결, 서비스데스크 운영 (담당: P5)
\item \textbf{가맹점 POS 단말 하드웨어 교체 및 그 비용} (P1 단말 예산 8.6억은 직영 477대 한정. 가맹 단말 교체 여부·비용 부담은 가맹 계약 개정 협상 결과에 따름)
\item 자사앱 자체의 UI/UX 개편 (담당: 앱 외주 벤더 계약, P1은 API 연동 제공까지)
\end{itemize}
\fi
\end{fieldbox}
\begin{fieldbox}
\textbf{경계 조건} \guide{(이 프로젝트는 ~의 기술 요건을 설계에 반영하되 ~자체는 하지 않는다)}\par
\ifwbblank
\foreach \i in {1,...,3}{\g{–} \rule{0pt}{8.5mm}\hrulefill\par}
\else
\begin{itemize}
\item P1은 P2 설비 인터페이스(15억, 별도 계약)의 데이터 규격을 P1 재고 원장 기준으로 제공하되, 설비측 구현과 시험 일정은 P2가 관리한다.
\item P1은 3PL WMS 준실시간 연동 명세를 작성하고 P1측 API를 구현하되, 대성로지스와의 계약 변경 협상은 물류본부(한동석) 소관이며 협상 미타결 시 야간 배치 연동을 폴백으로 유지한다.
\item P1은 규제 워크스트림(P4)이 확정한 아키텍처 요건을 설계에 반영하되, 요건 해석과 규제기관 대응은 P4·CPO 소관이다.
\item P1은 운영 이관 문서 12종 중 시스템 기술 문서를 작성하고, 운영 절차 문서는 P5가 작성한다.
\end{itemize}
\fi
\end{fieldbox}

% ------------------------------------------------------------------ 6
\secbar[70mm]{6. 주요 산출물 (Key Deliverables)}
\guide{시스템 구성요소뿐 아니라 관리 산출물(데이터 전환 결과, 인터페이스 설계문서, 소스·형상 인수, 운영 이관 문서)까지. 산출물마다 인수자(직책·실명)와 인수 시점. 소스·형상 인수를 빼놓는 것이 가장 흔한 감사 지적이다.}
{\footnotesize
\begin{tabularx}{\linewidth}{|L{36mm}|Y|L{38mm}|L{28mm}|}\hline
\hd{산출물} & \hd{설명} & \hd{인수자(직책·실명)} & \hd{인수 시점} \\ \hline
\ifwbblank
\rd{9mm}\g{[시스템 산출물]} & \g{[완료 상태를 한 줄로]} & & \g{[게이트/날짜]} \\ \hline
\rd{9mm} & & & \\ \hline
\rd{9mm} & & & \\ \hline
\rd{9mm} & & & \\ \hline
\rd{9mm} & & & \\ \hline
\rd{9mm}\g{[데이터 전환 결과]} & & & \\ \hline
\rd{9mm}\g{[설계문서 / 소스·형상 인수]} & & & \\ \hline
\rd{9mm}\g{[운영 이관 문서]} & & & \\ \hline
\else
신규 POS·매장 운영 시스템 & 615개 매장 배포 완료, 직영 단말 477대 설치 & 오정근 영업본부장 (현업 인수), 박세연 (기술 인수) & 컷오버 2027-02-28 \\ \hline
주문 허브 & 5개 채널 연동 완료, 옴니채널 시나리오 UAT 통과 & 오정근 & G3 2027-02-19 (UAT) \\ \hline
통합 재고 관리 & 3개 재고원(매장·센터·공장) 단일 원장, WMS 연동(준실시간 또는 폴백) & 한동석 물류본부장, 박세연 & G3 \\ \hline
가맹 발주 포털 & 388점 계정 발급, 리드타임 12h 설계 검증 & 오정근, 서정필(사용자 검수 참여) & G3 \\ \hline
통합 계층·데이터 전환 결과 & API GW·MDM·SAP 애드온 가동, 1,842 테이블 전환 검증 6종 통과 & 박세연 & 리허설 3회차 2027-01-29 → 컷오버 \\ \hline
인터페이스 설계문서 47본 & 문서 부재 13본 포함 전량 & 박세연 & G2 2026-11-27 \\ \hline
소스·형상 인수 & 발주사 저장소에 전 산출물 이관, 빌드 재현 확인 & 박세연, 발주사 형상관리 담당 & G4 2027-05-29 \\ \hline
운영 이관 기술 문서 & P5 SAC 대응 시스템 기술 문서 & 박준영 서비스운영팀장 & G4 \\ \hline
\fi
\end{tabularx}}

% ------------------------------------------------------------------ 7
\secbar[125mm]{7. 마일스톤 요약 (Summary Milestones)}
\guide{게이트와 외부 고정 날짜(규제, 계약, 다른 프로젝트 의존) 5~12개. 상세 일정이 아니다. \textbf{외부에서 고정된 날짜와 내부에서 정한 날짜를 구분}하고, 통과 조건과 판정자를 적는다 — 통과 조건이 없으면 마일스톤은 “지나가는 날짜”가 된다.}
{\footnotesize
\begin{tabularx}{\linewidth}{|L{22mm}|L{40mm}|L{22mm}|Y|L{26mm}|}\hline
\hd{날짜} & \hd{마일스톤} & \hd{외부 고정 여부} & \hd{통과 조건} & \hd{판정자} \\ \hline
\ifwbblank
\rd{8mm}\g{[YYYY-MM-DD]} & \g{[G0 착수 등]} & \g{[내부 / 외부(출처)]} & \g{[무엇이 충족되어야 통과인가]} & \g{[직책·실명]} \\ \hline
\rd{8mm} & & & & \\ \hline
\rd{8mm} & & & & \\ \hline
\rd{8mm} & & & & \\ \hline
\rd{8mm} & & & & \\ \hline
\rd{8mm} & & & & \\ \hline
\rd{8mm} & & & & \\ \hline
\rd{8mm} & & & & \\ \hline
\rd{8mm} & & & & \\ \hline
\rd{8mm} & & & & \\ \hline
\rd{8mm} & & & & \\ \hline
\rd{8mm} & \g{[종료 게이트 — 프로젝트 종료]} & & & \\ \hline
\else
2026-01-05 & G0 착수 & 내부 & 헌장 서명, 3역할 실명 지명 & 정미란 \\ \hline
2026-04-24 & P2 설비 발주 확정(T0) & \textbf{외부(P2)} & P1 재고 원장 인터페이스 규격 P2 전달 완료 & P2 PM \\ \hline
2026-04-30 & G1 기획 완료 & 내부 & 원가·일정 기준선, FP 기준선 확인서 서명, 조달 전략 확정 & 프로그램 이사회 \\ \hline
2026-08-14 & 가맹 계약 부속서 개정 & \textbf{외부(P3, 388점 동의)} & 가맹 단말 SW 배포 범위 확정 & 오정근 \\ \hline
2026-09-11 & 개인정보 보호법 개정 대응 완료 & \textbf{외부(규제)} & 처리항목 142개 설계 반영 검증 & 최영주 CPO \\ \hline
2026-11-27 & G2 중간 · 데이터 전환 정본 판정 & 내부 & 마스터 데이터 6종 확정 (이후 변경 시 610 테이블 재매핑 3.8억·6주) & 프로그램 이사회 \\ \hline
2027-01-29 & 데이터 전환 리허설 3회차 / 전자금융업 신청서 제출 & 내부 / \textbf{외부} & 전환 검증 6종 통과 & 박세연 / 최영주 \\ \hline
2027-02-12 & P2 기계적 준공(MC) & \textbf{외부(P2)} & 물류 인터페이스 통합시험 착수 가능 & P2 PM \\ \hline
2027-02-19 & G3 컷오버 승인 & 내부 & UAT 종료 기준 충족 + 롤백 절차 검증 & 프로그램 이사회 \\ \hline
2027-02-26~28 & 컷오버 주말 (금 20:00~일 06:00, 34시간) & \textbf{외부(성수기 창)} & 615개 매장 펌웨어 일괄배포, 롤백 시 2.1억·6영업일 & 정미란 (go/no-go) \\ \hline
2027-03-31 & 전자금융업 등록 심사 대응 & \textbf{외부(규제)} & P4 소관, P1 플랫폼 심사 대상 & 최영주 \\ \hline
2027-05-29 & \textbf{G4 운영 이관 완료 — P1 종료} & 내부 & SAC 미해소 0건 & 박준영 판정, 정미란 승인 \\ \hline
2028-03-31 & G5 편익 검토 & 내부 (P1 종료 후) & 연간 절감 판정 & 프로그램 이사회, 재무본부 \\ \hline
\fi
\end{tabularx}}

% ------------------------------------------------------------------ 8
\secbar[100mm]{8. 예산 요약 (Summary Budget)}
\guide{3층(승인 예산 / 집행 한도 / 계획 원가)과 각 층의 정의. 컨틴전시(PM 통제)와 관리예비비(스폰서 통제)를 반드시 구분한다. 사내 위임전결 한도와의 관계를 한 줄. “예산 168억”이라고만 적으면 PM과 CFO가 다른 숫자를 생각한다.}
{\footnotesize
\begin{tabularx}{\linewidth}{|L{38mm}|L{22mm}|Y|}\hline
\hd{층} & \hd{금액(억)} & \hd{정의 / 통제 주체} \\ \hline
\ifwbblank
\rd{8mm}승인 예산 & & \g{[누가 의결한 어떤 한도인가]} \\ \hline
\rd{8mm}집행 한도 & & \g{[승인 예산 − 무엇]} \\ \hline
\rd{8mm}계획 원가 & & \g{[집행 한도 − 무엇]} \\ \hline
\rd{8mm}컨틴전시 (PM 통제) & & \g{[비율, 용도, 집행 기록 방식]} \\ \hline
\rd{8mm}관리예비비 (스폰서 통제) & & \g{[유보분, 결재 주체]} \\ \hline
\rd{11mm}주요 내역 & \multicolumn{2}{Y|}{\g{[항목별 금액 — 비즈니스 케이스 요약 §5 인용]}} \\ \hline
\rd{11mm}사내 위임전결과의 관계 & \multicolumn{2}{Y|}{\g{[전결 한도, 결재 지침과 충돌 시 우선 규정]}} \\ \hline
\else
승인 예산 & 168.0 & 이사회 의결 계약·발주 한도 (P-08) \\ \hline
집행 한도 & 156.0 & 승인 예산 − 프로그램 관리예비비 12.0 (스폰서 결재 없이 쓸 수 없음) \\ \hline
계획 원가 & 145.8 & 집행 한도 − PM 통제 컨틴전시 10.2 \\ \hline
컨틴전시 (PM 통제) & 10.2 & 승인 예산의 6.07\%. 알려진 리스크 대응용. 집행 기록은 CCB에 월 보고 (2020 감사 지적 “컨틴전시 집행 승인 근거” 대응) \\ \hline
관리예비비 (스폰서 통제) & 12.0 & 프로그램 24억 중 P1 유보분. 스폰서 결재 \\ \hline
주요 내역 & \multicolumn{2}{Y|}{SI 68.2 / 상용SW 23.4 / 인프라 23.9 / 단말 8.6 / 데이터 전환 11.4 / 시험·감리 9.7 / PMO 12.6 / 컨틴전시 10.2 — 비즈니스 케이스 요약 §5} \\ \hline
사내 위임전결과의 관계 & \multicolumn{2}{Y|}{본부장 전결 5억. 단일 건 5,000만 원 이상 구매·용역·라이선스는 소관 본부장 결재 후 CFO 직접 결재 (2026-02-16 시행) — 재무본부 내부통제 지침} \\ \hline
\fi
\end{tabularx}}

% ------------------------------------------------------------------ 9
\secbar[\B{160mm}{70mm}]{9. 가정과 제약 (Assumptions \& Constraints)}
\guide{가정 = 참이라고 믿고 계획하지만 확인되지 않은 것(확인 시점·확인 책임자를 붙인다 → 나중에 리스크로 전환). 제약 = 바꿀 수 없는 것(출처 — 누가 정했는가 — 를 붙인다. 제약은 외부에서 온다). 둘을 섞지 말 것.}
\begin{fieldbox}
\textbf{가정} \guide{(~라고 가정한다. 확인 시점: [\ ], 확인 책임자: [\ ])}\par
\ifwbblank
\foreach \i in {1,...,6}{\g{–} \rule{0pt}{8.5mm}\hrulefill\ \g{(확인: \hspace{14mm}/ 책임: \hspace{14mm})}\par}
\else
\begin{itemize}
\item 대성로지스와의 3PL 계약 변경(준실시간 재고 연동)이 2026-06-30까지 타결된다. (확인: 한동석 물류본부장, 미타결 시 야간 배치 폴백)
\item 가맹 계약 부속서 개정이 2026-08-14까지 388점 서면 동의를 확보한다. (확인: 오정근, 미확보 시 가맹 SW 배포 범위를 동의 매장으로 축소)
\item 발주사 투입 인력 14명(대부분 겸임)의 P1 투입률이 평균 50\% 이상 유지된다. (확인: 박세연, 월 1회)
\item 자사앱 외주 벤더가 P1 API 규격에 맞춰 2026-12까지 앱 연동을 완료한다. (확인: 오정근·박세연, 앱 벤더 계약 변경 필요 시 P1 예산 외)
\item P2 기계적 준공이 2027-02-12에 달성되어 물류 인터페이스 통합시험이 G3 전에 가능하다. (확인: 프로그램 PMO, 지연 시 물류 부분 컷오버 분리 여부는 프로그램 이사회 결정)
\item 편익 산정 분모(순환실사 SKU 69\%)의 확장 방안이 M+3(2026-04)까지 확정된다. (확인: 한동석·나영선)
\end{itemize}
\fi
\end{fieldbox}
\begin{fieldbox}
\textbf{제약} \guide{(바꿀 수 없는 것. 출처: [\ ])}\par
\ifwbblank
\foreach \i in {1,...,7}{\g{–} \rule{0pt}{8.5mm}\hrulefill\ \g{(출처: \hspace{22mm})}\par}
\else
\begin{itemize}
\item 성수기(7~8월, 12월, 설·추석 전 2주)에는 매장 시스템 변경 불가. 컷오버 창은 2027-02-26~28로 고정. (출처: 영업본부, 프로그램 일정 기준선)
\item 규제 기한 2026-09-11, 2027-03-31은 협상 불가. P4가 22영업일 규제 버퍼를 두고 역산 관리. (출처: P4, CPO)
\item 컷오버 전까지 OD-POS 2.1(OS 지원 종료 상태)과 SAP ECC 6.0은 병행 운영. (출처: IT본부)
\item 자사앱 소스 형상관리는 외주 벤더 보유. (출처: 앱 유지보수 계약서)
\item 대성로지스 WMS는 벤더 소유, 커스터마이징 불가. (출처: 3PL 계약)
\item 본부장 전결 한도 5억, 5,000만 원 이상 CFO 결재. (출처: 위임전결규정, 내부통제 지침 2026-02-16)
\item ㈜온담푸드시스템은 별도 법인 — 공장 데이터 연동은 별도 계약·개인정보 처리 주체·보안 정책. (출처: 법무)
\item 수행사 계약: FP 12,400점 × 55만 원, 변경협의체 격주 수요일, 요구사항 베이스라인 월 1회 동결. (출처: P1 계약서 부속서 1)
\end{itemize}
\fi
\end{fieldbox}

% ------------------------------------------------------------------ 10
\secbar[\B{100mm}{60mm}]{10. 상위 리스크 (High-level Risks)}
\guide{스폰서가 서명 시점에 알아야 하는 리스크 상위 5~7개. “원인 → 사건 → 영향” 한 줄씩, 소유자 명시. 확률·영향 점수는 붙이지 않는다(Day2). PM 통제 밖의 의존(다른 프로젝트, 외부 계약, 규제)을 스폰서에게 알리는 자리다.}
\begin{fieldbox}
\ifwbblank
\foreach \i in {1,...,7}{\g{\i.} \rule{0pt}{9mm}\hrulefill\ \g{(소유자: \hspace{20mm})}\par}
\else
\begin{enumerate}
\item P2 설비 리드타임 34주→41주 통지 → MC 2027-02-12 미달 → 물류 인터페이스 통합시험이 G3 이후로 밀려 컷오버 물류 부분 분리 또는 창 이탈. (소유자: 프로그램 PMO장, P1측 대응: P1 PM)
\item 2022년과 같이 대성로지스가 데이터 개방·연동 방식 변경을 거부 → 준실시간 연동 불가 → 통합 재고의 센터 부분이 야간 배치로 남아 재고 정확도 목표 미달. (소유자: 한동석)
\item 가맹 388점 서면 동의 지연 또는 단말 비용 분쟁 → 가맹 SW 배포 범위 축소 → 가맹 판매 데이터 미수집, 폐기율·재고 정확도 편익의 가맹 부분 미실현. (소유자: 오정근)
\item CFO의 1차연도 집행 148억 축소 확정 → P1 2026년 배분 축소 → 라이선스 선약정 연기 또는 설계 인력 축소 → 컷오버 창 이탈. (소유자: 정미란, 대응안 제시: P1 PM)
\item OD-POS 원 개발자 1명(재직) 이탈 또는 병행 운영 중 보안 사고 → 인터페이스 13본 분석 불가·데이터 전환 지연 / 프로그램 신뢰 상실. (소유자: 박세연)
\item 발주사 겸임 인력 투입률 저하 → 발주사 미결정 대기 → 수행사 변경대가 청구(이재현 PM의 계약 협상 시 명시 입장). (소유자: 박세연, P1 PM)
\item 요구사항 동결 직전 주에 SRS 요구 집중 → 범위 허용범위 ±35건 초과 → FP 재산정 트리거. (소유자: P1 PM)
\end{enumerate}
\fi
\end{fieldbox}

% ------------------------------------------------------------------ 11
\secbar[80mm]{11. 주요 이해관계자 요약 (Key Stakeholders)}
\guide{의사결정 역할이 있는 사람과 성패에 결정적인 외부 이해관계자 8~12명 — 등록부(산출물 3)의 요약. 스폰서 칸에 부서·위원회명을 쓰지 말 것. 수행사 PM·가맹점주·노조·감사 같은 외부·비공식 이해관계자를 빼놓지 말 것.}
{\footnotesize
\begin{tabularx}{\linewidth}{|L{18mm}|L{34mm}|Y|Y|}\hline
\hd{성명} & \hd{직책} & \hd{프로젝트 역할} & \hd{핵심 관심사} \\ \hline
\ifwbblank
\rd{7.5mm} & & \g{[Executive]} & \\ \hline
\rd{7.5mm} & & \g{[Senior User]} & \\ \hline
\rd{7.5mm} & & \g{[Senior Supplier]} & \\ \hline
\rd{7.5mm} & & & \\ \hline
\rd{7.5mm} & & & \\ \hline
\rd{7.5mm} & & & \\ \hline
\rd{7.5mm} & & & \\ \hline
\rd{7.5mm} & & \g{[수행사 PM]} & \\ \hline
\rd{7.5mm} & & \g{[외부]} & \\ \hline
\rd{7.5mm} & & & \\ \hline
\rd{7.5mm} & & & \\ \hline
\rd{7.5mm} & & & \\ \hline
\else
정미란 & 재무본부장 CFO & Executive(스폰서), CCB 위원장, 관리예비비 결재 & 1차연도 집행 축소, 편익 근거 \\ \hline
오정근 & 영업본부장 & Senior User (P1·P3), 편익 실현 책임자 & 성수기 무중단, 가맹 반발 방지, 재계약률 91\% \\ \hline
박세연 & IT본부장 CIO & Senior Supplier, 아키텍처·인터페이스 최종 승인 & 자체 운영 가능, 벤더 비종속 \\ \hline
김선호 & 대표이사 & 프로그램 이사회 의장, 게이트 최종 승인 & IPO 전 가동, 일정 \\ \hline
한동석 & 물류본부장 & P2 Senior User, 3PL 계약 소관, 재고 정확도 책임 & 물류 자동화 우선, 처리량 목표 근거 \\ \hline
최영주 & 준법지원팀장 CPO & CCB 위원, 규제 요건 해석 & 시행일 준수, 검토 의견 제공(결정 아님) \\ \hline
박준영 & 서비스운영팀장 & P5 인수자, SAC 판정권 & 운영 문서·SLA 없이는 인수 불가 \\ \hline
나영선 & 품질안전본부장 & CCB 위원, 품질 데이터 요구 제공 & 검사기록 디지털화 \\ \hline
이재현 & ㈜한빛디지털 PM & 수행사 PM, 변경협의체 대표 & FP 기준선, 변경대가, 손익률 6.5\% \\ \hline
서정필 & 가맹점주협의회장 & 외부, 가맹 포털 실사용자 대표, 부속서 개정 협상 상대 & 데이터 소유, 단말 비용, 가맹 매출 \\ \hline
윤태호 & 내부감사팀장 & 관리 체계 서면 질의자 & 편익 근거, 컨틴전시 승인 근거 \\ \hline
김미르 & 노조 물류지부장 & 이천센터 246명 대표 & 처리량 목표와 인력 \\ \hline
\fi
\end{tabularx}}

% ------------------------------------------------------------------ 12
\secbar[\B{150mm}{90mm}]{12. 거버넌스와 PM 권한 (Governance \& PM Authority)}
\guide{의사결정 구조(누가 무엇을 승인하는가), PM이 스폰서 보고 없이 결정할 수 있는 허용범위(tolerance), 초과 \textbf{예상} 시 에스컬레이션 경로와 시한, 사내 위임전결과의 충돌 시 우선 규정. 헌장에서 가장 협상이 필요한 칸이다.}
\begin{fieldbox}
\textbf{의사결정 구조} \guide{(스폰서 / 운영위 / CCB / 변경협의체 — 각각 무엇을 결정하는가)}\par
\ifwbblank
\lines{4}{9mm}
\else
프로그램 이사회(게이트 판정, 관리예비비 24억, 프로그램 재정의 — 분기 1회 + 게이트) → 프로그램 스폰서 정미란(P1 예산 계층 간 이동, 허용범위 초과 승인, 관리예비비 12억) → 프로그램 운영위(월 1회, 프로젝트 간 조정) → P1 PM. 변경은 CCB(7인, 월 1회 셋째 목요일, 정족수 5, 반대의견 실명 기재)가 사내 승인, 변경협의체(발주 3 + 수행 3, 격주 수요일 14:00)가 계약상 변경대가·FP 재산정 트리거 판정. 기준선 재확인은 분기 1회 또는 FP 누적 변동 ±3\% 초과 시.
\fi
\end{fieldbox}
\needspace{32mm}\vspace{1.5mm}
\textbf{\small PM 허용범위 (프로젝트 계층)}\par
{\footnotesize
\begin{tabularx}{\linewidth}{|Y|Y|Y|Y|Y|Y|}\hline
\hd{원가} & \hd{일정} & \hd{범위} & \hd{품질} & \hd{리스크} & \hd{편익} \\ \hline
\ifwbblank
\rd{10mm}\g{[+n\% (금액)]} & \g{[+n영업일]} & \g{[±n\% (건수)]} & \g{[결함밀도 등]} & \g{[EMV n억]} & \g{[−n\%p]} \\ \hline
\else
+6\% (10.08억) & +15영업일 & ±3\% (±35건) & 결함밀도 0.4건/FP & EMV 8억 & −5\%p \\ \hline
\fi
\end{tabularx}}
\begin{fieldbox}
\textbf{“PM은 위 허용범위 안에서 스폰서 보고 없이 결정한다.”}\par
\ifwbblank
\g{에스컬레이션: 허용범위 초과 예상 시 [PM → 스폰서 (n영업일 내) → 프로그램 이사회]} \lines{2}{9mm}
\g{사내 전결 규정과의 관계:} \lines{2}{9mm}
\else
P1 PM은 위 허용범위 안에서 스폰서 보고 없이 결정한다. 단, 사내 위임전결규정에 따라 단일 건 5억 초과는 허용범위 내라도 CFO 결재를 병행하고, 5,000만 원 이상 발주는 내부통제 지침(2026-02-16)에 따른다. 허용범위 초과가 \textbf{예상}되는 시점(초과한 시점이 아니라)에 PM은 5영업일 내 스폰서에게 예외 보고(원인·영향·대안 3개·권고)를 제출하고, 스폰서는 10영업일 내 결정하거나 프로그램 이사회에 부의한다. 작업패키지 계층 허용범위(원가 +2\%, 일정 +5영업일)는 수행사 PM에게 위임한다. (회사 정본 §5.6 초안 기준)
\fi
\end{fieldbox}

% ------------------------------------------------------------------ 13
\secbar[50mm]{13. 우선순위 (Priority among Constraints)}
\guide{일정·원가·범위·품질이 충돌할 때 무엇을 지키고 무엇을 양보하는지 — \textbf{스폰서가} 결정해 적는다. 고정 / 조정 가능 / 수용 세 칸과 이유 한 문장. 모두 “고정”이면 우선순위가 아니라 희망이다.}
{\footnotesize
\begin{tabularx}{\linewidth}{|Y|Y|Y|Y|}\hline
\hd{고정 (fixed)} & \hd{조정 가능 (flexible)} & \hd{수용 (accept)} & \hd{결정자 / 이유} \\ \hline
\ifwbblank
\rd{26mm}\g{[지켜야 하는 축과 그 값]} & \g{[양보할 축과 범위]} & \g{[허용범위 내 감수할 것]} & \g{[스폰서 직책·실명, 날짜, 이유 한 문장]} \\ \hline
\else
컷오버 창 2027-02-26~28 (일정) / 집행 한도 156억 (원가) / 규제 요건·가용성 99.5\% (품질 필수 항목) & 범위 — SRS 570건 중 옴니채널 고도화(매장 간 재고 이동, 픽업 예약)와 BI 리포트는 컷오버 이후 릴리스로 이연 가능 & 원가 허용범위 +6\% 내 컨틴전시 사용, 일정 +15영업일 & 정미란 스폰서, 2026-01-05. 이유: 컷오버 창을 놓치면 다음 창은 2027년 하반기이며 IPO 예비심사와 겹친다. 원가는 부채비율 목표 때문에 집행 한도를 넘길 수 없다. 따라서 조정 변수는 범위뿐이다. \\ \hline
\fi
\end{tabularx}}

% ------------------------------------------------------------------ 14
\secbar[\B{110mm}{60mm}]{14. 완료 기준 (Completion / Exit Criteria)}
\guide{정상 종료 조건 3~5줄(“오픈”이 아니라 “운영 이관 완료”인가, 판정자·승인자) + 조기 종료(중단) 조건 2~3줄(부의자, 손실 확정 절차). 중단 조건이 없으면 중단 결정을 아무도 내리지 못한다.}
\begin{fieldbox}
\textbf{정상 종료} \guide{(판정: [\ ], 승인: [\ ])}\par
\ifwbblank
\lines{4}{9mm}
\else
G4(2027-05-29) 시점에 ① SAC 미해소 항목 0건 또는 해소 기한·책임자가 서명된 목록, ② 소스·형상 발주사 인수 100\%, ③ 인터페이스 47본 설계문서 인수, ④ 컷오버 후 3개월 가용성 99.5\% 달성, ⑤ 종료 보고서·교훈 기록 제출. 판정: 박준영 서비스운영팀장(SAC), 박세연(형상·문서). 승인: 정미란.
\fi
\end{fieldbox}
\begin{fieldbox}
\textbf{조기 종료(중단) 조건} \guide{(부의: [\ ], 손실 확정: [\ ])}\par
\ifwbblank
\lines{3}{9mm}
\else
① G1 또는 G2에서 편익 실현률 전망이 손익분기 66\% 미만으로 판정될 때, ② 컷오버 창을 2회 연속 이탈(2027-02, 2027년 하반기)할 때, ③ 프로그램 이사회가 프로그램을 재정의할 때. 중단 부의: 정미란 → 프로그램 이사회. 손실 확정: 재경팀장이 확정 손실(기집행액, 해지 위약금, 라이선스 잔여 의무)을 30일 내 산정하고 내부감사팀에 통보(2020 감사 지적 “손실 확정 절차 부재” 대응).
\fi
\end{fieldbox}

% ------------------------------------------------------------------ 15
\secbar[75mm]{15. 서명 (Sign-off)}
\guide{Executive(스폰서), Senior User, Senior Supplier, PM의 실명·직책·날짜·서명. 프로그램 하위 프로젝트면 프로그램 PMO장의 확인. Senior User(현업)의 서명이 없으면 나중에 “우리는 동의한 적 없다”가 나온다.}
{\footnotesize
\begin{tabularx}{\linewidth}{|L{34mm}|Y|L{26mm}|L{42mm}|}\hline
\hd{역할} & \hd{직책·실명} & \hd{날짜} & \hd{서명} \\ \hline
\rd{11mm}Executive (스폰서) & \B{}{재무본부장 CFO 정미란} & \B{}{2026-01-05} & \B{}{(서명)} \\ \hline
\rd{11mm}Senior User & \B{}{영업본부장 오정근} & \B{}{2026-01-05} & \B{}{(서명)} \\ \hline
\rd{11mm}Senior Supplier & \B{}{IT본부장 CIO 박세연} & \B{}{2026-01-05} & \B{}{(서명)} \\ \hline
\rd{11mm}프로젝트 매니저 & \B{}{P1 PM} & \B{}{2026-01-05} & \B{}{(서명)} \\ \hline
\rd{11mm}확인 (프로그램 PMO장) & \B{}{프로그램 PMO장} & \B{}{2026-01-05} & \B{}{(서명)} \\ \hline
\end{tabularx}}

\end{document}
